\section{Experimental Setup}
\label{sec:methodology}

\subsection{Tasks, Data and Scoring}
We use two tasks with short answers that can be checked automatically. \textbf{GSM8K}~\cite{cobbe2021gsm8k} consists of grade-school math word problems whose answer is a single number. \textbf{HotpotQA}~\cite{yang2018hotpotqa} consists of multi-hop questions over Wikipedia. In its distractor setting, each question comes with two gold paragraphs and eight distractor paragraphs, and all ten paragraphs are given to the model. The data are 500 questions from the GSM8K test set and 500 from the HotpotQA validation set, sampled with seed 42. Each dataset is randomly split (seed 42) into a development split $D_{\text{dev}}$ with 20\% of the questions, used only to choose thresholds, and a test split $D_{\text{test}}$ with the remaining 80\%, used once for all reported results.

Every prompt that produces an answer asks the model to finish with a line of the form \texttt{Final answer: <answer>}, and only this span is scored. For GSM8K we extract the number, remove formatting such as commas, currency signs and trailing zeros, and compare it with the gold number. For HotpotQA we apply the official answer normalisation (lower-casing and removing punctuation, articles and extra whitespace) and compute exact match (EM) and token-level F1. The correctness label $Y$ is EM on the extracted answer. To measure how much apparent harm is caused only by changes in answer form, we also use a \emph{lenient match}, which counts an answer as correct if the normalised gold answer is contained in the normalised extracted answer.

\subsection{Models and Inference}
We use the open instruction-tuned model \modelA{}~\cite{qwen2025qwen25} with its own chat template, served with vLLM 0.6.6.post1 on two NVIDIA T4 GPUs (16\,GB each) in 16-bit floating point. The initial answer, the revision and the IoE baseline use greedy decoding, so the only randomness in the pipeline comes from the self-consistency samples, each of which is drawn with its own random seed.

\subsection{Pipeline and Confidence Signals}
For each question we first generate the initial answer $A_0$. We then compute three confidence scores $S(A_0)\in[0,1]$, each in a separate call so that one signal cannot influence another. The full prompts are listed in Appendix~\ref{app:prompts}.
\begin{enumerate}
    \item \textbf{Verbalised confidence $S_{\text{verb}}$.} In a new turn after $A_0$, the model is asked to state how confident it is that its final answer is correct, as a whole number from 0 to 100, which we divide by 100. If the reply cannot be parsed, we set $S_{\text{verb}}=1$, so that the answer is never revised by this gate, and we report the number of such cases.
    \item \textbf{Self-consistency agreement $S_{\text{sc}}$.} We draw $N=10$ additional samples with the same prompt as $A_0$, at temperature $0.7$ and top-$p$ $0.95$. With $k$ the number of samples whose normalised extracted answer equals that of $A_0$, we set $S_{\text{sc}}=k/N$. We also report the variant with $N=5$, which uses the first five samples.
    \item \textbf{Self-verification $S_{\text{verif}}$ (P(True)).} Following Kadavath et al.~\cite{kadavath2022language}, the model sees the question (for HotpotQA also the ten paragraphs) together with the extracted answer of $A_0$, and must choose between ``(A) True'' and ``(B) False''. With $p_A$ and $p_B$ the next-token probabilities of the two options, summed over their tokenisation variants, we set $S_{\text{verif}}=p_A/(p_A+p_B)$. The gold answer never appears in this prompt.
\end{enumerate}
Finally, the model revises $A_0$ in the same conversation with a fixed critique-and-revise prompt, producing $A_1$. The revision is computed for every question, whatever its confidence. This lets us evaluate all gating policies offline on exactly the same outputs, so differences between policies come only from the decision of whether to revise.

Table~\ref{tab:signals} summarises what each signal needs. The signals differ in two practical ways: whether they require access to token probabilities, and whether their cost comes mainly from reading the input (prefill) or from generating text (completion).

\begin{table}[h!]
\caption{The three confidence signals and what they require per question, in addition to generating $A_0$.}
\centering
\small
\begin{tabular}{lcccc}
\toprule
Signal & Extra calls & Token probabilities & Main cost & Output \\
\midrule
$S_{\text{verb}}$  & 1  & not needed & short completion & integer 0--100 \\
$S_{\text{sc}}$    & 10 & not needed & 10 full completions & agreement $k/N$ \\
$S_{\text{verif}}$ & 1  & needed     & prefill (question and context) & $p_A/(p_A+p_B)$ \\
\bottomrule
\end{tabular}
\label{tab:signals}
\end{table}

\subsection{Gating, Baselines and Threshold Selection}
A gate with signal $S$ and threshold $\tau$ outputs $A_1$ if $S(A_0)<\tau$ and $A_0$ otherwise. We evaluate gates on $S_{\text{verb}}$, $S_{\text{sc}}$ and $S_{\text{verif}}$, and on the simple combination $S_{\text{avg}}=\tfrac12(S_{\text{verb}}+S_{\text{sc}})$, which, like its two parts, needs no access to token probabilities. For each gate, the threshold $\tau^*$ is chosen on $D_{\text{dev}}$ among all observed score values as the one with the highest dev accuracy; ties are broken in favour of the threshold that revises fewer questions. The chosen $\tau^*$ is then applied once to $D_{\text{test}}$.

We compare the gates against five baselines:
\begin{itemize}
    \item \textbf{Never revise:} output $A_0$.
    \item \textbf{Always revise:} output $A_1$.
    \item \textbf{Majority vote:} output the most frequent extracted answer among $A_0$ and the ten self-consistency samples, a strong baseline that uses the same samples as $S_{\text{sc}}$.
    \item \textbf{IoE}~\cite{li2024confidence}: a single revision prompt that tells the model to keep its answer if it is confident and to change it otherwise, using the prompt of the original paper.
    \item \textbf{Oracle gate:} revise exactly the questions for which $A_0$ is wrong. It uses the gold labels and serves only as an upper bound, not as a practical method.
\end{itemize}

\paragraph{Majority vote as a revision step.} To test the analytical model with a revision step that does repair errors, we also use the majority answer over the ten self-consistency samples, denoted $R_{\text{mv}}$, in place of $A_1$, and gate it with $S_{\text{sc}}$ and $S_{\text{avg}}$ in exactly the same way. This needs no additional model calls, because the samples are already computed for $S_{\text{sc}}$.

\subsection{Metrics and Cost Model}
\textbf{Transitions.} Each question falls into one of four cells: $W\to R$ ($Y_0=0$, $Y_1=1$), $R\to W$ ($Y_0=1$, $Y_1=0$), $R\to R$ or $W\to W$. For a gated method, only the questions it revises can change cell. From the always-revise transitions we estimate $a$, $f$, $b$ and $\kappa$ as defined in Section~\ref{sec:theory}.

\textbf{Discrimination.} For each signal we report the AUROC for predicting $Y_0=1$ on $D_{\text{test}}$, with 95\% confidence intervals from 1{,}000 bootstrap resamples of the test questions.

\textbf{Significance.} Two methods are evaluated on the same questions, so we compare their accuracies with the exact (binomial) McNemar test on the discordant pairs, against both never revising and always revising. We report exact $p$-values rather than significance stars.

\textbf{Cost.} For question $i$ and a gated method we count
\begin{equation}
C_i \;=\; C_i^{A_0} \;+\; C_i^{S} \;+\; \mathbb{1}[S_i<\tau]\;C_i^{\text{rev}},
\label{eq:cost}
\end{equation}
where $C_i^{A_0}$, $C_i^{S}$ and $C_i^{\text{rev}}$ are the tokens used for the initial answer, for computing the signal and for the revision. We keep prompt (prefill) tokens and generated (completion) tokens separate. Prefill matters here because the verification and revision calls re-read the question and, for HotpotQA, all ten paragraphs. For majority voting the cost is the initial answer plus the ten samples.
